\documentclass[11pt]{article}
% Source-PDF: 英文资料 ENGLISH MATERIALS/Wikipedia_link-cut_tree.pdf
\usepackage{oipl-vn}

\title{Cây link/cut}
\author{}
\date{}

\begin{document}
\maketitle

\origtitle{Link/cut tree}
\sourcepath{english/Wikipedia\_link-cut\_tree.pdf}

\section*{Tổng quan}

\term{Cây link/cut} là một cấu trúc dữ liệu dùng để biểu diễn một rừng, tức
một tập các cây có gốc. Cấu trúc này hỗ trợ các thao tác sau:
\begin{itemize}
  \item thêm vào rừng một cây chỉ gồm một đỉnh;
  \item cho một đỉnh trong một cây, tách đỉnh đó cùng toàn bộ cây con của nó
  khỏi cây hiện tại;
  \item gắn một đỉnh làm con của một đỉnh khác;
  \item cho một đỉnh, tìm gốc của cây chứa nó. Khi thực hiện thao tác này trên
  hai đỉnh khác nhau, ta có thể kiểm tra chúng có thuộc cùng một cây hay không.
\end{itemize}

Rừng được biểu diễn có thể gồm các cây rất sâu. Nếu chỉ lưu rừng bằng các cây
con trỏ cha thông thường, việc tìm gốc của một đỉnh có thể mất nhiều thời gian.
Ngược lại, nếu mỗi cây trong rừng được biểu diễn bằng cây link/cut, ta có thể
xác định một phần tử thuộc cây nào trong thời gian khấu hao $O(\log n)$.
Hơn nữa, ta có thể nhanh chóng cập nhật tập các cây link/cut khi rừng biểu
diễn thay đổi; cụ thể, việc hợp nhất bằng thao tác \texttt{link} và tách bằng
thao tác \texttt{cut} đều có thời gian khấu hao $O(\log n)$.

\begin{center}
\begin{tabular}{lcc}
\hline
\textbf{Thao tác} & \textbf{Trung bình} & \textbf{Trường hợp xấu nhất} \\
\hline
\texttt{Link} & $O(\log n)$ khấu hao & $O(\log n)$ \\
\texttt{Cut} & $O(\log n)$ khấu hao & $O(\log n)$ \\
\texttt{Path} & $O(\log n)$ khấu hao & $O(\log n)$ \\
\texttt{FindRoot} & $O(\log n)$ khấu hao & $O(\log n)$ \\
\hline
\end{tabular}
\end{center}

Cây link/cut chia mỗi cây trong rừng được biểu diễn thành các đường đi rời nhau
theo đỉnh. Mỗi đường đi được biểu diễn bằng một cấu trúc dữ liệu phụ trợ,
thường là cây splay; bài báo gốc ra đời trước cây splay nên dùng cây tìm kiếm
nhị phân thiên lệch. Các đỉnh trong cấu trúc phụ trợ được sắp theo độ sâu của
chúng trong cây được biểu diễn tương ứng. Trong biến thể \textit{phân hoạch
ngây thơ}, các đường đi được xác định bởi các đường đi và đỉnh được truy cập
gần nhất, tương tự cây Tango. Trong \textit{phân hoạch theo kích thước}, đường
đi được xác định theo con nặng nhất, tức con có nhiều hậu duệ nhất, của đỉnh
đang xét. Cách này làm cấu trúc phức tạp hơn, nhưng giảm chi phí thao tác từ
$O(\log n)$ khấu hao xuống $O(\log n)$ trong trường hợp xấu nhất. Cấu trúc có
ứng dụng trong nhiều bài toán luồng mạng và trên các tập dữ liệu động.

Trong công trình gốc, Sleator và Tarjan gọi cây link/cut là ``cây động'', hoặc
``dynamic dyno trees''.

\section{Cấu trúc}

Ta xét một cây mà mỗi đỉnh có số con tùy ý và các con không có thứ tự, rồi tách
cây đó thành các đường đi. Cây ban đầu này được gọi là \term{cây được biểu
diễn}. Các đường đi được biểu diễn bên trong bằng các cây phụ trợ; ở đây ta
dùng cây splay. Trong cây phụ trợ, thứ tự từ trái sang phải biểu diễn đường đi
từ gốc đến đỉnh cuối của đường đi đó. Những đỉnh kề nhau trong cây được biểu
diễn nhưng không nằm trên cùng một đường đi ưu tiên, và do đó không nằm trong
cùng một cây phụ trợ, được nối bằng một \term{con trỏ cha của đường}. Con trỏ
này được lưu ở gốc của cây phụ trợ biểu diễn đường đi.

\subsection{Đường đi ưu tiên}

Khi thực hiện truy cập đến một đỉnh $v$ trong cây được biểu diễn, đường đi được
đi qua trở thành \term{đường đi ưu tiên}. \term{Con ưu tiên} của một đỉnh là
đứa con cuối cùng nằm trên đường truy cập, hoặc là rỗng nếu lần truy cập cuối
cùng là đến chính $v$, hoặc nếu chưa từng truy cập nhánh cây cụ thể đó.
\term{Cạnh ưu tiên} là cạnh nối con ưu tiên với đỉnh đang xét.

Trong một phiên bản khác, các đường đi ưu tiên được xác định bởi con nặng nhất.

\section{Thao tác}

Các thao tác ta quan tâm là \texttt{FindRoot(Node v)}, \texttt{Cut(Node v)},
\texttt{Link(Node v, Node w)} và \texttt{Path(Node v)}. Mọi thao tác đều được
cài đặt bằng thủ tục con \texttt{Access(Node v)}. Khi ta truy cập một đỉnh
$v$, đường đi ưu tiên của cây được biểu diễn được đổi thành đường đi từ gốc $R$
của cây được biểu diễn đến đỉnh $v$. Nếu một đỉnh trên đường truy cập trước đó
có con ưu tiên $u$, còn đường đi mới đi qua con $w$, thì cạnh ưu tiên cũ bị xóa
và được đổi thành con trỏ cha của đường, còn đường đi mới sẽ đi qua $w$.

\subsection{\texttt{Access}}

Sau khi thực hiện \texttt{Access} tại đỉnh $v$, đỉnh này không còn con ưu tiên
nào và nằm ở cuối đường đi. Vì các đỉnh trong cây phụ trợ được khóa theo độ
sâu, mọi đỉnh ở bên phải $v$ trong cây phụ trợ phải bị tách ra. Với cây splay,
thao tác này khá đơn giản: ta splay tại $v$, đưa $v$ lên gốc của cây phụ trợ.
Sau đó ta tách cây con phải của $v$; đó chính là mọi đỉnh nằm dưới $v$ trên
đường đi ưu tiên cũ. Gốc của cây bị tách sẽ có một con trỏ cha của đường, và
ta cho con trỏ đó trỏ đến $v$.

Tiếp theo ta đi ngược lên cây được biểu diễn đến gốc $R$, phá và đặt lại đường
đi ưu tiên khi cần. Để làm điều này, ta đi theo con trỏ cha của đường từ $v$;
vì $v$ hiện là gốc của cây phụ trợ, ta truy cập trực tiếp được con trỏ đó. Nếu
đường đi chứa $v$ đã chứa gốc $R$ thì, do các đỉnh được khóa theo độ sâu, $R$
là đỉnh trái nhất trong cây phụ trợ; khi đó con trỏ cha của đường là rỗng và
thao tác truy cập kết thúc.

Ngược lại, con trỏ dẫn đến một đỉnh $w$ trên một đường đi khác. Ta muốn phá
đường đi ưu tiên cũ của $w$ rồi nối nó với đường đi chứa $v$. Để làm vậy, ta
splay tại $w$ và tách cây con phải của nó, đặt con trỏ cha của đường của cây
bị tách thành $w$. Vì mọi đỉnh được khóa theo độ sâu, và mọi đỉnh trên đường
đi chứa $v$ đều sâu hơn mọi đỉnh trên đường đi chứa $w$, ta chỉ cần nối cây
của $v$ làm con phải của $w$. Sau đó ta splay lại tại $v$; do $v$ là con của
gốc $w$, thao tác này chỉ xoay $v$ lên gốc. Ta lặp toàn bộ quá trình cho đến
khi con trỏ cha của đường của $v$ là rỗng; khi đó $v$ nằm trên cùng đường đi
ưu tiên với gốc $R$ của cây được biểu diễn.

\subsection{\texttt{FindRoot}}

\texttt{FindRoot} là thao tác tìm gốc của cây được biểu diễn chứa đỉnh $v$.
Vì thủ tục \texttt{Access} đưa $v$ vào đường đi ưu tiên, trước hết ta thực hiện
\texttt{Access}. Khi đó $v$ nằm trên cùng đường đi ưu tiên, và do đó trong
cùng cây phụ trợ, với gốc $R$. Vì các cây phụ trợ được khóa theo độ sâu, gốc
$R$ là đỉnh trái nhất của cây phụ trợ. Ta chỉ cần liên tục đi sang con trái
của $v$ cho đến khi không thể đi tiếp; đỉnh thu được là gốc $R$. Gốc này có
thể nằm ở độ sâu tuyến tính, là trường hợp xấu nhất đối với cây splay, nên ta
splay nó để lần truy cập sau nhanh hơn.

\subsection{\texttt{Cut}}

Ở đây ta muốn cắt cây được biểu diễn tại đỉnh $v$. Trước hết ta truy cập $v$.
Việc này đặt mọi phần tử nằm thấp hơn $v$ trong cây được biểu diễn vào cây con
phải của $v$ trong cây phụ trợ. Mọi phần tử hiện nằm trong cây con trái của
$v$ là các đỉnh cao hơn $v$ trong cây được biểu diễn. Do đó ta tách con trái
của $v$; cây này vẫn giữ liên kết với cây được biểu diễn ban đầu thông qua con
trỏ cha của đường. Lúc này $v$ là gốc của một cây được biểu diễn. Việc truy cập
$v$ cũng phá đường đi ưu tiên bên dưới $v$, nhưng cây con đó vẫn giữ liên kết
với $v$ thông qua con trỏ cha của đường.

\subsection{\texttt{Link}}

Nếu $v$ là gốc của một cây và $w$ là một đỉnh trong cây khác, thao tác
\texttt{Link} nối hai cây chứa $v$ và $w$ bằng cách thêm cạnh $(v,w)$, làm cho
$w$ trở thành cha của $v$. Để thực hiện, ta truy cập cả $v$ và $w$ trong hai
cây tương ứng rồi đặt $w$ làm con trái của $v$. Vì $v$ là gốc, còn các đỉnh
trong cây phụ trợ được khóa theo độ sâu, sau khi truy cập $v$ thì $v$ không có
con trái trong cây phụ trợ, vì nó có độ sâu nhỏ nhất. Thêm $w$ làm con trái
trên thực tế làm cho $w$ trở thành cha của $v$ trong cây được biểu diễn.

\subsection{\texttt{Path}}

Trong thao tác này, ta muốn tính một hàm gộp nào đó trên tất cả các đỉnh hoặc
cạnh thuộc đường đi từ gốc $R$ đến đỉnh $v$, chẳng hạn \texttt{sum},
\texttt{min}, \texttt{max}, \texttt{increase}, v.v. Ta thực hiện
\texttt{Access} tại $v$, qua đó nhận được một cây phụ trợ chứa toàn bộ các
đỉnh trên đường đi từ gốc $R$ đến $v$. Cấu trúc dữ liệu có thể được mở rộng để
lưu thông tin cần truy vấn, chẳng hạn giá trị nhỏ nhất, giá trị lớn nhất, hoặc
tổng chi phí trong cây con; thông tin trên một đường đi khi đó có thể được trả
về trong thời gian hằng số sau khi đã thực hiện truy cập.

\section{Phân tích}

\texttt{Cut} và \texttt{Link} có chi phí $O(1)$ cộng với chi phí của
\texttt{Access}. \texttt{FindRoot} có cận trên khấu hao $O(\log n)$ cộng với
chi phí của \texttt{Access}. Tùy cách cài đặt, cấu trúc dữ liệu có thể được mở
rộng bằng các thông tin bổ sung, chẳng hạn đỉnh có giá trị nhỏ nhất hoặc lớn
nhất trong các cây con, hoặc tổng. Do đó \texttt{Path} có thể trả về thông tin
này trong thời gian hằng số cộng với cận của \texttt{Access}.

Vì vậy, để tìm thời gian chạy, phần còn lại là chặn chi phí của
\texttt{Access}. Thao tác \texttt{Access} dùng splay, mà ta biết có cận trên
khấu hao $O(\log n)$. Phần phân tích còn lại xét số lần cần splay. Số này bằng
số lần thay đổi con ưu tiên, tức số cạnh thay đổi trong đường đi ưu tiên, khi
ta đi ngược lên cây.

Ta chặn \texttt{Access} bằng kỹ thuật \term{phân rã nặng-nhẹ}.

\subsection{Phân rã nặng-nhẹ}

Kỹ thuật này gọi một cạnh là nặng hoặc nhẹ tùy theo số đỉnh trong cây con. Ký
hiệu $\operatorname{size}(v)$ là số đỉnh trong cây con của $v$ trong cây được
biểu diễn. Một cạnh được gọi là \term{nặng} nếu
\[
  \operatorname{size}(v) > \frac{1}{2}
  \operatorname{size}(\operatorname{parent}(v)).
\]
Do đó mỗi đỉnh có nhiều nhất một cạnh nặng. Cạnh không nặng được gọi là
\term{nhẹ}.

\term{Độ sâu nhẹ} là số cạnh nhẹ trên đường đi từ gốc đến một đỉnh $v$. Ta có
\[
  \text{độ sâu nhẹ} \le \lg n,
\]
vì mỗi lần đi qua một cạnh nhẹ, số đỉnh giảm ít nhất một nửa so với cây con của
cha.

Vậy một cạnh trong cây được biểu diễn có thể thuộc một trong bốn loại:
nặng-ưu tiên, nặng-không ưu tiên, nhẹ-ưu tiên, hoặc nhẹ-không ưu tiên. Trước
hết ta chứng minh cận trên $O(\log^2 n)$.

\subsubsection{\texorpdfstring{Cận trên $O(\log^2 n)$}{Cận trên O(log binh phuong n)}}

Mỗi thao tác splay trong \texttt{Access} đóng góp hệ số $\log n$, nên để chứng
minh cận trên $O(\log^2 n)$, ta cần chặn số lần truy cập bởi $\log n$.

Mỗi lần cạnh ưu tiên thay đổi, một cạnh ưu tiên mới được tạo ra. Vì vậy ta đếm
số cạnh ưu tiên được tạo. Trên một đường đi bất kỳ có nhiều nhất $\log n$ cạnh
nhẹ, nên có nhiều nhất $\log n$ cạnh nhẹ chuyển thành ưu tiên.

Số cạnh nặng chuyển thành ưu tiên có thể là $O(n)$ trong một thao tác cụ thể,
nhưng chỉ là $O(\log n)$ theo nghĩa khấu hao. Trên một dãy thao tác, ban đầu có
thể có $n-1$ cạnh nặng chuyển thành ưu tiên, vì toàn bộ cây được biểu diễn có
nhiều nhất $n-1$ cạnh nặng. Từ đó trở đi, số cạnh nặng chuyển thành ưu tiên
bằng số cạnh nặng đã chuyển thành không ưu tiên ở một bước trước. Với mỗi cạnh
nặng chuyển thành không ưu tiên, phải có một cạnh nhẹ chuyển thành ưu tiên. Ta
đã thấy số cạnh nhẹ có thể chuyển thành ưu tiên nhiều nhất là $\log n$. Vì vậy
trong $m$ thao tác, số cạnh nặng chuyển thành ưu tiên là
\[
  m\log n + (n-1).
\]
Với đủ nhiều thao tác, cụ thể $m > n-1$, giá trị trung bình là $O(\log n)$.

\subsubsection{\texorpdfstring{Cải thiện thành cận trên $O(\log n)$}{Cải thiện thành cận trên O(log n)}}

Ta đã chặn số lần thay đổi con ưu tiên bởi $O(\log n)$. Nếu chứng minh được
mỗi lần thay đổi con ưu tiên có chi phí khấu hao $O(1)$, ta sẽ chặn được thao
tác \texttt{Access} bởi $O(\log n)$. Điều này được thực hiện bằng phương pháp
thế năng.

Gọi $s(v)$ là số đỉnh nằm dưới $v$ trong cây của các cây phụ trợ. Khi đó hàm
thế năng là
\[
  \Phi = \sum_v \log s(v).
\]
Ta biết chi phí khấu hao của thao tác splay được chặn bởi
\[
  \operatorname{cost}(\operatorname{splay}(v))
  \le 3\bigl(\log s(\operatorname{root}(v))-\log s(v)\bigr)+1.
\]

Sau khi splay, $v$ là con của đỉnh $w$ được trỏ bởi con trỏ cha của đường, nên
ta có
\[
  s(v) \le s(w).
\]
Sử dụng bất đẳng thức này cùng chi phí khấu hao của truy cập, ta thu được một
tổng dạng lồng nhau được chặn bởi
\[
  3\bigl(\log s(R)-\log s(v)\bigr)
  + O(\text{số lần thay đổi con ưu tiên}),
\]
trong đó $R$ là gốc của cây được biểu diễn. Vì số lần thay đổi con ưu tiên là
$O(\log n)$ và $s(R)=n$, ta thu được cận khấu hao $O(\log n)$.

\section{Ứng dụng}

Cây link/cut có thể được dùng để giải bài toán kết nối động trên đồ thị không
chu trình. Với hai đỉnh $x$ và $y$, chúng liên thông khi và chỉ khi
\[
  \texttt{FindRoot}(x) = \texttt{FindRoot}(y).
\]
Một cấu trúc dữ liệu khác cũng dùng được cho mục đích này là cây Euler tour.

Trong bài toán luồng cực đại, cây link/cut có thể cải thiện thời gian chạy của
thuật toán Dinic từ $O(V^2E)$ xuống $O(VE\log V)$.

\section{Xem thêm}

\begin{itemize}
  \item Cây splay.
  \item Phương pháp thế năng.
\end{itemize}

\section{Đọc thêm}

\begin{sloppypar}
\begin{itemize}
  \item Sleator, D. D.; Tarjan, R. E. (1983). ``A Data Structure for Dynamic
  Trees''. \textit{Proceedings of the thirteenth annual ACM symposium on Theory
  of computing -- STOC '81}, p. 114. doi:\url{https://doi.org/10.1145/800076.802464}.

  \item Sleator, D. D.; Tarjan, R. E. (1985). ``Self-Adjusting Binary Search
  Trees''. \textit{Journal of the ACM}, 32(3): 652.
  doi:\url{https://doi.org/10.1145/3828.3835}.

  \item Goldberg, A. V.; Tarjan, R. E. (1989). ``Finding minimum-cost
  circulations by canceling negative cycles''. \textit{Journal of the ACM},
  36(4): 873. doi:\url{https://doi.org/10.1145/76359.76368}. Ứng dụng cho tuần
  hoàn chi phí nhỏ nhất.

  \item ``Link-Cut trees'', ghi chú bài giảng về cấu trúc dữ liệu nâng cao,
  Spring 2012, bài giảng 19. Giáo sư Erik Demaine; người ghi chép: Justin
  Holmgren (2012), Jing Jian (2012), Maksim Stepanenko (2012), Mashhood Ishaque
  (2007).

  \item \url{http://compgeom.cs.uiuc.edu/~jeffe/teaching/datastructures/2006/notes/07-linkcut.pdf}
\end{itemize}
\end{sloppypar}

\end{document}
